\documentclass[10pt,a4paper]{amsart}
\usepackage[margin=19mm]{geometry}
\usepackage{amsmath,amssymb,mathtools}
\usepackage[scr=rsfs,cal=euler]{mathalfa}
\usepackage{xfrac}
\usepackage[unicode,hidelinks,bookmarks=false]{hyperref}
\newcommand{\ie}{i.e.\ }
\newcommand{\cf}{cf.\ }
\newcommand{\resp}{respectively\ }
\newcommand{\Subsection}{\subsection}
\providecommand{\nobreakdash}{\nobreak\hbox{-}}
\setlength{\emergencystretch}{2em}
\multlinegap0pt
\usepackage{bm}
\usepackage{bbm}
\usepackage{dsfont}
\usepackage{xspace}
\usepackage{cancel}
\usepackage{mathtools}
\usepackage{amsthm}
\makeatletter
\@ifundefined{theorem}{\newtheorem{theorem}{Theorem}}{}
\@ifundefined{ex}{\newtheorem{ex}[theorem]{Example}}{}
\@ifundefined{prop}{\newtheorem{prop}[theorem]{Proposition}}{}
\makeatother
\title[Enumerating binary words restricted by subsequence frequency]{Enumerating binary words restricted by subsequence frequency}
\author{Glenn Bruda}
\date{Source: arXiv:2607.02578v1 (CC BY 4.0)}
\begin{document}
\maketitle

\noindent\emph{Source and licence.} Enumerating binary words restricted by subsequence frequency. Version arXiv:2607.02578v1 (CC BY 4.0), distributed under the Creative Commons Attribution 4.0 International licence, as recorded in the arXiv licence metadata for this exact version and shown on the abstract page https://arxiv.org/abs/2607.02578v1. Full rights notice: licenses/arxiv-2607.02578-CC-BY-4.0.txt.\par\smallskip

\noindent\textit{Editorial scope.} Counted unit: the Prop whose source label is \texttt{combinatorial\_formula\_for\_number\_of\_occurrences} in the article (source0.tex lines 695--700, source label \texttt{combinatorial\_formula\_for\_number\_of\_occurrences}), delivered together with the article's own complete original proof of it (source0.tex lines 710--740, one \texttt{proof} environment of 31 lines). No mathematics is written, repaired or re-derived: statement and proof are the article's own text line for line. The proof is detached from the statement in the source: the article states the result at lines 695--700 and prints its proof at lines 710--740, and the proof environment's own header names the statement, so the association is the article's own. Required context retained uncounted: inverse\_of\_algorithm (lines 630--640). Every definition, display and label of those lines is retained unchanged. Omitted from this package and not redistributed: 1--629, 641--694, 701--709, 741--1128. Nothing inside a retained excerpt is omitted or altered: every retained body is delivered exactly as the article's lines are written, so no omission receipt of any kind accompanies this package. Citations: the retained excerpts carry no citation command at all, so no bibliography entry of the article is transcribed and no reference is redistributed. Reference adaptation: none was needed and none was made; every reference inside the delivered unit resolves inside it, so no unresolved reference or citation is produced. Printed slips kept literal: no printed word, symbol, hypothesis or number of the counted statement or of its proof was altered. Layout and class: the article's own class is \texttt{amsart} with packages including graphicx, appendix, comment, enumerate, esvect, bm, amsmath, amssymb, mathtools, mathrsfs, color, ulem, url, bbm; the wrapper is the lane's pinned \texttt{amsart} editorial wrapper at 10pt on A4, which restates the theorem-like environments the retained text needs and adds the article's own macro definitions that the retained text needs (none), with extra wrapper packages (bm, bbm, dsfont, xspace, cancel, mathtools). The delivered numbering is the unit's own. Assets: no retained span contains a figure environment, a graphics-inclusion command, a PDF-page inclusion, a table or a TikZ picture; no image file is copied into this package, the package declares no asset dependency and no image file is redistributed with it. Licence: version arXiv:2607.02578v1 of this article is distributed under the Creative Commons Attribution 4.0 International licence, as recorded in the arXiv licence metadata for this exact version and shown on the abstract page https://arxiv.org/abs/2607.02578v1. A full rights notice is delivered at licenses/arxiv-2607.02578-CC-BY-4.0.txt. Characters: every retained excerpt is pure ASCII, so the delivered engine, XeLaTeX with its default Latin Modern text font, reports no missing character.
\begin{ex}\label{inverse_of_algorithm}
    \normalfont Let $p=111001111$ and $w=10011110101001101110101$. Let $\mathbf{i}$ be the occurrence of $p$ described by the underlined letters in $\underline{1}00\underline{11}1101\underline{0}10\underline{011}0\underline{11}10101$. Begin by placing a bar in $\mathbf{i}$ immediately after the last letter of each run of $p$, obtaining
    \begin{align*}
        \underline{1}00\underline{11}\Big|1101\underline{0}10\underline{0}\Big|\underline{11}0\underline{11}\Big|10101.
    \end{align*}
    Now, slide each bar to the right until it does not lie in the middle of a single run of $w$. If a bar already does not lie in a single run of $w$ (e.g., the second bar above), then do not slide that bar. Applying this procedure to the above yields
    \begin{align*}
        \underline{1}00\underline{11}11\Big|01\underline{0}10\underline{0}\Big|\underline{11}0\underline{11}1\Big|0101,
    \end{align*}
    from which we obtain the partition $1001111,010100,110111$ of the prefix $1001111010100110111$ of $w$ into $r=3$ intervals. Finally, the corresponding tuple of $\mathbf{i}$ is then the number of runs in each interval of this partition, returning $(3,5,3)\in \mathscr{C}_{3,\text{\normalfont odd}}^{\leq 15}$. 
\end{ex}

\begin{prop}\label{combinatorial_formula_for_number_of_occurrences}
    Let $p=x_1^{u_1}\cdots x_r^{u_r}$ and $w=y_1^{v_1}\cdots y_{R}^{v_R}$ be binary words. Let $R'$ be the number of runs of $w$ which are not NBRs and set $\epsilon=1-\delta_{x_1,y_1}$. Then
    \begin{align}\label{eqn_comb_formula}
        c_p(w)=\sum_{C\in\mathscr{C}_{r,\text{\normalfont odd}}^{\leq R'}}\prod_{A_i\in C([R'])}\left({\sum_{j\in A_i}v_{j+\epsilon}\choose u_i}-{\sum_{j\in A_i\setminus\{\max(A_i)\}}v_{j+\epsilon}\choose u_i}\right).
    \end{align}
\end{prop}

\begin{proof}[Proof of Proposition~\ref{combinatorial_formula_for_number_of_occurrences}]
    We first note that the proof reduces to the case where $w$ has no NBRs, i.e., $x_1=y_1$ and $x_r=y_R$. Indeed, letting $w'=y_{1+\epsilon}^{v_{1+\epsilon}}\cdots y_{R'+\epsilon}^{v_{R'+\epsilon}}$, it follows that $c_p(w)=c_p(w')$ and $w'$ has no NBRs.
    
    Thus, restricting to the case where $x_1=y_1$ and $x_r=y_R$, let $\mathbf{i}=(i_1<\cdots<i_{\ell})\in\mathfrak{C}_p(w)$. For each $1\leq a\leq r$, let $s(a)=\sum_{j=1}^{a}u_j$ and $I_a=i_{s(a)}$ and $g_{\mathbf{i}}(a)$ be such that $w_{I_a}$ lies in the run $y_{g_{\mathbf{i}}(a)}^{v_{g_{\mathbf{i}}(a)}}$. In words, $g_{\mathbf{i}}(a)$ is the index of the run in $w$ which contains the last letter of the $a$\textsuperscript{th} run of $p$ with respect to $\mathbf{i}$. In Example~\ref{inverse_of_algorithm}, we have $g_{\mathbf{i}}(1)=3$, $g_{\mathbf{i}}(2)=3+5=8$, and $g_{\mathbf{i}}(3)=3+5+3=11$. 
    
    Define $C=(g_{\mathbf{i}}(1),g_{\mathbf{i}}(2)-g_{\mathbf{i}}(1),g_{\mathbf{i}}(3)-g_{\mathbf{i}}(2),\dots,g_{\mathbf{i}}(r)-g_{\mathbf{i}}(r-1))$ to be the $r$-tuple associated to $\mathbf{i}$. We claim that each element of $C$ is odd. As $x_1=y_1$, an occurrence of an odd-indexed run of $p$ in $w$ must lie only in the odd-indexed runs of $w$. Similarly, an occurrence of an even-indexed run of $p$ in $w$ must lie only in the even-indexed runs of $w$. So $g_{\mathbf{i}}(n)$ is of the same parity as $n$. Thus, each element of $C$ is odd. Observing that the sum of the elements of $C$ telescopes to $g_{\mathbf{i}}(r)\leq R$, we have
    \begin{align*}
        \{r\text{-tuple~associated~to~}\mathbf{i}:\mathbf{i}\in\mathfrak{C}_{p}(w)\}\subseteq\mathscr{C}_{r,\text{odd}}^{\leq R}.
    \end{align*}
    So
    \begin{align}\label{partition_of_occurrences}
        c_p(w)=\sum_{C\in\mathscr{C}_{r,\text{odd}}^{\leq R}}\left|\{\mathbf{i}\in\mathfrak{C}_{p}(w):C=r\text{-tuple~associated~to~}\mathbf{i}\}\right|.
    \end{align}
    To complete the proof, we now enumerate the summands in \eqref{partition_of_occurrences}. Let $(c_1,\dots,c_r)=C\in \mathscr{C}_{r,\text{odd}}^{\leq R}$ and $\mathbf{i}\in\mathfrak{C}_{p}(w)$ be such that $C$ equals the $r\text{-tuple associated to }\mathbf{i}$. Let $1\leq t\leq r$. For notational convenience, set $g_{\mathbf{i}}(0)=0$. 
    
    %By construction, the last letter of the run $x_t^{u_t}$ occurs in the run $y_{g_{\mathbf{i}}(t)}^{v_{g_{\mathbf{i}}(t)}}$. Furthermore, the first letter of the run $x_t^{u_t}$ occurs in a run after the run $y_{g_{\mathbf{i}}(t-1)}^{v_{g_{\mathbf{i}}(t-1)}}$. 
    
    By construction, the last letter of the $t$\textsuperscript{th} run of $p$ occurs in the $g_{\mathbf{i}}(t)$\textsuperscript{th} run of $w$. Furthermore, the first letter of the $t$\textsuperscript{th} run of $p$ occurs in a run after the $g_{\mathbf{i}}(t-1)$\textsuperscript{th} run of $w$. Thus, as $x_t=y_b$ only if $b$ has the same parity as $t$, the $t$\textsuperscript{th} run of $p$ may only occur in runs with indices in $\{g_{\mathbf{i}}(t-1)+1,g_{\mathbf{i}}(t-1)+3\dots,g_{\mathbf{i}}(t)-2,g_{\mathbf{i}}(t)\}$, with at least one letter occurring in the $g_{\mathbf{i}}(t)$\textsuperscript{th} run of $w$. Noting that $\sum_{h=1}^{T}c_h=g_{\mathbf{i}}(T)$, the occurrences of the $t$\textsuperscript{th} run of $p$ with the restriction of $C$ are enumerated by the $u_t$-element subsets of
    \begin{align}\label{big_set_prop_proof}
        \left\{1,\dots,v_{g_{\mathbf{i}}(t-1)+1}+v_{g_{\mathbf{i}}(t-1)+3}+\cdots+v_{g_{\mathbf{i}}(t)-2}+v_{g_{\mathbf{i}}(t)}\right\}=\left\{1,\dots,\sum_{j\in A_t} v_j\right\}
    \end{align}
    containing at least one of the $v_{g_{\mathbf{i}}(t)}=v_{\max(A_t)}$ largest elements of \eqref{big_set_prop_proof}. The number of these subsets is
    \begin{align*}
        {\sum_{j\in A_t}v_{j}\choose u_t}-{\sum_{j\in A_t\setminus\{\max(A_t)\}}v_{j}\choose u_t},
    \end{align*}
    which may be seen via an all-minus-bad strategy. Thus, the size of $\{\mathbf{i}\in\mathfrak{C}_{p}(w):C=r\text{-tuple~associated~to~}\mathbf{i}\}$ is
    \begin{align}\label{complicated_summand}
        \prod_{A_t\in C([R])}\left({\sum_{j\in A_t}v_{j}\choose u_t}-{\sum_{j\in A_t\setminus\{\max(A_t)\}}v_{j}\choose u_t}\right),
    \end{align}
    yielding \eqref{eqn_comb_formula}.
\end{proof}

\end{document}
